\documentclass[12pt]{article}
\usepackage[utf8]{inputenc}
\usepackage[T1]{fontenc}
\usepackage{hyperref}
\title{Southeast Pashai Language Resources for K-12 Curriculum Development}
\author{Compiled by Claude Sonnet 5.5 (claude-sonnet-5-5) for Andr\'as Kornai}
\date{}
\begin{document}
\maketitle

\section{Introduction}
Southeast Pashai (ISO 639-3: psi; Glottocode: sout2672) is an Indo-European language of Afghanistan with roughly 180,000 speakers. Web searches found very little material aimed at schools. No primer, textbook or dictionary of the language itself turned up; the resources below are mostly reference or general-profile pages. This report is deliberately short rather than padded.

\section{Primary Educational Resources}
No school primers, textbooks or curriculum documents for Southeast Pashai were found in the search results. Joshua Project describes the language as written and as having a Bible translation that has been started, but this is not an educational resource in itself.

\section{Linguistic and Reference Materials}
\begin{itemize}
\item Pashayi languages (Wikipedia overview): \url{https://en.wikipedia.org/wiki/Pashayi_languages}
\item Lamuwal and Baker (2013), IPA illustration of Southeastern Pashayi, Journal of the International Phonetic Association (PDF): \url{https://core-cms.cambridgecore.org/core/services/aop-cambridge-core/content/view/8CDDFBEC9998A428ACA1FFB023555626/S0025100313000133a.pdf/southeastern-pashayi.pdf}
\item PHOIBLE phoneme inventory for Southeast Pashai: \url{https://phoible.org/languages/sout2672}
\end{itemize}
The search summaries also mention Morgenstierne's grammar of Pashayi and Bashir (2003), and orthography work by the Pashai Language Committee; no direct URLs for these were seen, so none is given.

\section{Religious and Community Resources}
\begin{itemize}
\item Joshua Project profile, Southeast Pashayi in Afghanistan: \url{https://joshuaproject.net/people_groups/15025}
\item Joshua Project language page for psi: \url{https://www.joshuaproject.net/index.php/languages/psi}
\end{itemize}

\section{Digital and Audio Resources}
\begin{itemize}
\item Global Recordings Network language page for psi (the search summary reports no audio recordings available): \url{https://www.globalrecordings.net/en/language/psi}
\end{itemize}

\section{Curriculum Development Context}
The language is spoken as a first language by its community, but no evidence of schooling in it was found. A curriculum would have to begin with orthography standardization and new materials, building on the phonological description and the Pashai Language Committee's earlier alphabet work, with community leadership.

\end{document}